\section{Conclusions}
\label{sec:conclusions}

We compared towing kites, rigid wing sails and Flettner rotors on thirteen classes of cargo ship, six Atlantic routes and three fuels, with each rig sized for the lowest cost per tonne-kilometre on one routed model. The claim is conditional: if each rig reaches the goals of \cref{tab:goals}, the choice comes out as follows.

Use a kite on container ships, rotors on car carriers, and either on bulk carriers, tankers and general cargo ships, chosen by route. On the container ships the kite saves about \tbd{11}\,\% of the cost per tonne-kilometre on e-methanol and leads on every route except the Megamax's to Dakar; on fuel oil it saves about \tbd{3}\,\% of the cost and \tbd{22}\,\% of the CO$_2$, and the rotors about \tbd{0 to 2}\,\%. On the car carrier the rotors save \tbd{12.5}\,\% on e-methanol and lead on every route; on fuel oil they save \tbd{4}\,\% of the cost and \tbd{23}\,\% of the CO$_2$. On the ships that fill their deadweight kites and rotors both save about \tbd{13}\,\% on e-methanol, level within the accuracy of the model, so either works: rotors lead on the routes to Dakar and Tubarão, with much wind from abeam and ahead, kites on the routes to Guadeloupe, Houston and Bilbao, and the kite wherever the deck takes only the few rotors makers offer today. On fuel oil the kite leads on nearly all of them, saving \tbd{4.4}\,\% of the cost and \tbd{25}\,\% of the CO$_2$, and like every best rig on fuel oil it pays back in years rather than months (\cref{sec:industry}). On liquid hydrogen the order is the same on the container ships and the car carrier, while on the ships that fill their deadweight the rotors lead by about a point. Wing sails save about half as much as kites and rotors and lead nowhere: a square metre of rotor costs about three times a square metre of wing and pulls about five times as hard (\cref{sec:res_forces}). An e-methanol ship with its best rig still costs about \tbd{40}\,\% more than the same ship on fuel oil without a rig, against about \tbd{60}\,\% without its rig, so the best rig closes about a third of the premium and no rig makes the zero-emission ship as cheap as today's ship on fuel oil.

The split follows what each rig pays in cargo (\cref{sec:res_groups}). On the deadweight-limited bulk carriers and tankers the rotors pull harder per square metre, but their mass costs them \tbd{0.4 to 1.2} points of cargo and the kite's less than \tbd{0.2}, which leaves the two within \tbd{0.8} points on e-methanol. On the space-limited feeder, Panamax and Neo-Panamax container ships the slots the rotors block cost them \tbd{5 to 7} points and the kite's mast less than half a point; without that cost the rotors would lead. The car carrier's rig costs no cargo at all, so there the rotors' harder pull decides. The line between the groups is close for the largest container ships, and a box \tbd{7}\,\% lighter or \tbd{9}\,\% heavier than the mean moves the Megamax or the Neo-Panamax across it, which widens the kite's lead to \tbd{6} points on the one and puts the rotors ahead on the other. On deadweight-limited ships the deck itself costs no cargo, as real installations confirm, but it limits how many deck units fit: held to the units makers offer today, the rotors fall \tbd{4 to 8} points behind the kite on each of the six bulk carriers and tankers for which a maker gives one.

Part of every saving is time: since cost sets the speed, a ship with a kite sails a median of \tbd{6.3}\,\% faster and one with rotors \tbd{10.8}\,\%, and a ship held to its schedule takes the gain as fuel instead. The published view that kites suit slow ships rests on a kite of one size; a larger kite takes back the share of fuel it saves at a higher speed (\cref{sec:res_speed}).

Three inputs can change which rig to choose (\cref{sec:res_situations}). \tbd{The fuel price decides most: on the Kamsarmax the kite leads up to about 14~USD/GJ of fuel oil and the rotors above it, while on the \SI{140}{m} ship the kite leads over the whole range. A cargo value above about 6000~USD/t lets the rotors pass the kite on the \SI{140}{m} ship, and a high discount rate weighs on the deck rigs, whose cost is mostly capital, more than on the kite. No change to the wing sails alone puts them ahead.}

The goals are not equally far off, and those of the kite decide the order. Rotors have about eight years of service behind their figures and wing sails about two, while no kite has a service record. Reaching the canopy life goal rather than \num{2000} flight hours is worth \tbd{1.3 to 2.6} points of saving and reaching the canopy price goal \tbd{0.8 to 1.7}, and the kite at its cheapest size can take only \tbd{0.16 to 3.3} crashes a year before it falls behind the next rig (\cref{sec:res_missed}), so reliability decides whether a ship kite is worth building at all. A router limited to a 3-day forecast keeps about \tbd{70}\,\% of the wing sails' saving; \tbd{the kite and the rotors are pending the forecast run}.

A VTOL kite, a rigid wing that takes off and lands vertically on its own propellers, flies at ship scale but does not pay at any size, even with every input at the bound that favours it, at about four times the soft kite's yearly cost per unit of thrust (\cref{sec:res_rigid}).

These results hold under series-production prices that no kite maker has reached, deck rigs of up to 23 units where no ship carries more than five, kites of up to \num{16000}~m$^2$, sixteen times the largest designed, and one wind profile for all regions (\cref{sec:discussion}). On the largest ships the deck, a physical limit, sets the size of the wings and rotors, while the kite stops at the bound of the search, a limit of the model, so there its saving is a lower bound.
